% New experiment version: use with the paired tables in tables/legacy_tables.tex.

This is a draft based exclusively on the new measured results in this package. Renumber tables and figures after inserting it. Update the method description to the actual two-detour encoding before using this text.

\subsection{Experimental protocol}

We evaluated rerouting on 45 legacy synthetic traffic instances, comprising 15 cases on each of three complete meshes with 4, 5, and 6 nodes. We additionally constructed 72 controlled cases varying off-peak demand, congestion-trigger boundaries, localized and distributed bottlenecks, single-edge capacity loss, and network-wide capacity loss, together with two sparse-network route-availability diagnostics. All input graphs, loads, capacities, packet allocations, candidate routes and QUBOs were saved before solver execution. Within each traffic instance, QAOA and BPSO used the same decision variables, objective coefficients and decoding rules.

Each fixed traffic input was evaluated using five solver seeds. QAOA used Qiskit's StatevectorSampler at p=1 with 1024 shots per sampler evaluation and COBYLA with a maximum of 50 iterations. BPSO used 30 particles, 20 iterations, w=0.5 and c1=c2=1.5. Exhaustive enumeration established two distinct references: the minimum QUBO objective J and the minimum attainable peak load/capacity ratio over the represented route assignments. Objective optimality gaps were computed only against the former. Every enumerated QUBO cost was independently checked against the objective calculated from decoded traffic loads. The encoding allocates each generated 150 veh/hr packet to one of two detours, under a global 10-variable budget; retaining a packet on its original edge is not one of those binary choices.

\subsection{Scenario outcomes}

Across the 45 legacy traffic inputs, the QAOA mean peak improved the unmanaged state in 40 cases, was unchanged in 1, and worsened in 4. Comparing five-seed mean peaks, QAOA was lower than BPSO in 0 cases, tied in 45, and was higher in 0. Median relative peak reductions across scenario-level QAOA means were 4 nodes: 11.57\%, 5 nodes: 16.05\%, 6 nodes: 18.71\%. These are descriptive results across five solver seeds per fixed traffic state, with sample standard deviations reported alongside paired per-scenario values. We do not interpret absent differences as statistical equivalence or claim computational advantage.

The controlled experiments isolate demand severity and spare route capacity from capacity-loss effects. For lane-closure and weather experiments, input demand was held fixed while the specified capacities were reduced. Off-peak and trigger-boundary cases test whether intervention is activated appropriately; sparse-network cases identify the current requirement for two candidate detours. All deteriorating and no-action cases remain in the tables and outcome denominators.

All nine controlled single-bottleneck cases improved peak occupancy, with relative reductions of 28.21\% to 34.38\%. All nine separated-bottleneck cases also improved (10.42\% to 19.74\%). In contrast, distributed-peak cases worsened in 7/9 cases and were unchanged in 2/9; no case improved. All six reduced-capacity weather cases worsened; the three unreduced-capacity controls required no action. These outcomes were shared by BPSO. They support a conditional interpretation: this rerouting model helps localized demand when useful detours exist, but can amplify congestion under widespread load/capacity stress. The attribution is limited to these declared synthetic inputs.

\subsection{Interpretation of marginal and negative outcomes}

Positive relative peak relief does not necessarily clear congestion. We therefore report absolute relief in percentage points, residual edge load/capacity ratios, secondary congestion and capacity exceedances. In 13 legacy cases, the lowest peak among exact J minima exceeded the minimum peak attainable over candidate assignments, demonstrating a mismatch between the optimized weighted balancing objective and the reported peak metric. In 3 cases, every represented assignment worsened the initial peak, identifying a limitation of the mandatory-detour candidate set. Solver approximation is assessed separately through the exact J gap.

Exact-only diagnostics on four difficult legacy cases vary the packet budget and allocation order while holding traffic fixed. These diagnostics avoid attributing deterioration to the qubit budget merely because its limit is reached. They also expose that a larger packet budget changes which traffic is forced to detour and need not monotonically improve peak occupancy. These static synthetic experiments establish behavior within the represented decision problem; they do not establish real-world travel-time improvements, hardware speedup or dynamic queue recovery.
